%!TEX root = main.tex

\subsection{Overview: \GNC Algorithm with Non-minimal Solvers} % Alternation of }
% \subsection{Overview: GNC Algorithm with Non-minimal Solvers}
\label{sec:GNC-algorithm-overview}

We start by providing an overview of the proposed algorithm, then we delve into technical details and 
tailor the approach to two specific robust cost functions (Section~\ref{sec:GNC-GM-TLS}).  
Instead of optimizing directly\revise{~the robust estimation problem}~\eqref{eq:generalEstimation}, we use \GNC and, 
at each \emph{outer} iteration, we fix a $\mu$ and optimize: 
%
\bea \label{eq:GNCestimation}
\min_{\vxx \in \calX} \sumAllPointsi \rho_\mu( r(\vy_i,\vxx) ).
\eea 
%
Since non-minimal solvers cannot be used directly to solve~\eqref{eq:GNCestimation}, 
%the surrogate function is difficult to solve globally 
we use the Black-Rangarajan duality and rewrite~\eqref{eq:GNCestimation} using the corresponding outlier process:
%
\bea \label{eq:GNC-outlierProcess}
\min_{\vxx \in \calX, w_i \in [0,1]} \sumAllPointsi  \left[ w_i r^2(\vy_i,\vxx) + \Phi_{\rho_\mu}(w_i) \right].
\eea 
%
As discussed in Section~\ref{sec:GNC-GM-TLS} and~\cite{Black96ijcv-unification}, it is easy to compute the penalty terms $\Phi_{\rho_\mu}(\cdot)$ even for the surrogate function $\rho_\mu(\cdot)$. 

Finally, we solve~\eqref{eq:GNC-outlierProcess} by alternating optimization, where  
 at each \emph{inner} iteration we first optimize over the $\vxx$ (with fixed $w_i$), and 
 then we optimize over $w_i$, $i=1,\ldots,N$ (with fixed $\vxx$).  % which admits a closed-form solution.
 In particular, at inner iteration $t$, we perform the following:
 %we alternate the following update steps:
 %each iteration alternates the following update steps:

\begin{enumerate}
\item {\bf Variable update}: minimize~\eqref{eq:GNC-outlierProcess} with respect to $\vxx$ with fixed weights $w\at{t-1}_i$:
% Solve the weighted least squares optimization with fixed weights $w_i$:
\bea \label{eq:variableUpdate}
\vxx\at{t} = \argmin_{\vxx \in \calX} \sumAllPointsi w\at{t-1}_i r^2(\vy_i, \vxx),
\eea
where we dropped the second term in~\eqref{eq:GNC-outlierProcess} which does not depend on $\vxx$. 
Problem~\eqref{eq:variableUpdate} is simply a weighted version of the outlier-free problem~\eqref{eq:leastSquares}, 
 hence it can be solved globally using certifiably optimal non-minimal solvers.  
% and obtain the estimated residuals $\hat{r}_i = r(\vy_i, \vxx^\star )$. This step can be solved with certifiable optimality using global solvers.

\item {\bf Weight update}: minimize~\eqref{eq:GNC-outlierProcess} with respect to $w_i$ ($i=1,\ldots,N$) 
with fixed $\vxx\at{t}$:
% Update the weights with estimated residuals $\hat{r}_i$:
\bea \label{eq:weightUpdate}
\vw\at{t} = \argmin_{w_i \in [0,1]} \sumAllPointsi \left[ w_i r^2(\vy_i, \vxx\at{t})  + \Phi_{\rho_\mu}(w_i)\right],
\eea
where $r^2(\vy_i, \vxx\at{t})$ is a constant for fixed $\vxx$, and the expression of $\Phi_{\rho_\mu}(\cdot)$ depends on the choice of robust cost function. As we discuss in Section~\ref{sec:GNC-GM-TLS}, the weight update can be typically solved in closed form. 
\end{enumerate}

 % The pseudo-code of the \GNC algorithm is given in Algorithm~\ref{}.
 % \LC{Hank: please add pseudocode of algorithm \vspace{5cm}}
The process is then repeated for changing values of $\mu$, where each change of $\mu$ increases the amount of non-convexity. 
% The algorithm terminates when \LC{....} 
% \LC{something that remains unclear is how many inner iterations we need..}
\setcounter{theorem}{1}
\begin{remark}[Teaching an Old Dog New Tricks]
While the combination of \GNC and Black-Rangarajan duality has been investigated in related works~\cite{Black96ijcv-unification,Zhou16eccv-fastGlobalRegistration}, its applicability has been limited by the lack of global solvers for the variable update in~\eqref{eq:variableUpdate}. 
For instance,~\cite{Zhou16eccv-fastGlobalRegistration} focuses on a specific problem (point cloud registration) where~\eqref{eq:variableUpdate} can be solved in closed form~\cite{Arun87pami,Horn87josa},  
%using Horn's or Arun's methods~\cite{Arun87pami,Horn87josa}, 
while~\cite{Black96ijcv-unification} focuses on a Markov Random Field formulation for which global solvers and heuristics exist~\cite{Hu19ral-fuses}.
 One of the main insights behind our approach is that modern non-minimal solvers (developed over the last 5 years) 
 allow solving~\eqref{eq:variableUpdate} globally for a broader class of problems,  
 % largely expand the domain of application of \GNC, 
 including spatial perception problems such as SLAM, mesh registration, and object localization from images.
 %which can be effectively used in complex 
% and uses a heuristic approach. 
\end{remark}

\isExtended{
\begin{remark}[Generality]
The algorithm described in this section
%~Section~\ref{sec:GNC-algorithm-overview} 
is applicable to any estimation problem and robust 
cost function as long as: (i) we can develop a surrogate function $\rho_\mu$ for \GNC and find an analytic penalty function $\Phi_{\rho_\mu}$, % to obtain the outlier process in eq.~\eqref{eq:GNC-outlierProcess},  
(ii) we have an efficient way to solve the variable update~\eqref{eq:variableUpdate}, and 
(iii) we have an efficient way to solve the weight update~\eqref{eq:weightUpdate}.
Condition (i) is relatively mild and~\cite{Black96ijcv-unification} already provides analytic expressions for the most common robust costs. Condition (iii) is also mild since the optimization splits into $N$ scalar optimizations, whose solution can be typically computed in closed form. 
Therefore, the only ``stringent'' requirement is the availability of non-minimal solver for~\eqref{eq:variableUpdate}. 
Luckily, non-minimal solvers exist for several key spatial perception applications as the ones considered in Section~\ref{sec:applications}. 
% a global solver for the variable update~\eqref{eq:variableUpdate}: however, we have already observed that the variable update has the same form of the outlier free 
\end{remark}
}{}

% \red{
% \begin{remark}[Generality]
% The algorithm described in this section
% %~Section~\ref{sec:GNC-algorithm-overview} 
% is applicable to any estimation problem and robust 
% cost function as long as: (i) we can develop a surrogate function $\rho_\mu$ for \GNC, 
% (ii) we can find an analytic penalty function $\Phi_{\rho_\mu}$ to obtain the outlier process in eq.~\eqref{eq:GNC-outlierProcess},  
% (iii) we have an efficient way to solve the variable update~\eqref{eq:variableUpdate}, and 
% (iv) we have an efficient way to solve the weight update~\eqref{eq:weightUpdate}.
% Conditions (i) and (ii) are relatively mild and~\cite{Black96ijcv-unification} already provides analytic expressions for the most common robust costs. Condition (iv) is also mild since the optimization splits into $N$ scalar optimizations, whose solution can be typically computed in closed form. 
% Therefore, the only ``stringent'' requirement is the availability of non-minimal solver for~\eqref{eq:variableUpdate}. 
% Luckily, non-minimal solvers exist for several key spatial perception applications as the ones considered in Section~\ref{sec:applications}. 
% % a global solver for the variable update~\eqref{eq:variableUpdate}: however, we have already observed that the variable update has the same form of the outlier free 
% \end{remark}
% }
% With the outlier process $\Phi(w_i)$, Black-Rangarajan duality ensures problem~\eqref{eq:optiOutlierProcess} has the same solution as the original problem~\eqref{eq:generalEstimation}, but explicitly allows for solving the optimization by alternatively updating the unknown variable $\vxx$ and the weight $w_i$ for each measurement.
% \begin{proposition}[Primal-dual Solution for Outlier Process] \label{prop:primalDualOutlierProcess}
% A local minimum of the robust line process optimization problem~\eqref{eq:optiOutlierProcess} can be found in two steps:
% \begin{enumerate}
% \item {\bf Primal}: Solve the weighted least squares optimization with fixed weights $w_i$:
% \bea \label{eq:variableUpdate}
% \vxx^\star = \argmin_{\vxx \in \calX} \sumAllPointsi w_i r^2(\vy_i, \vxx) 
% \eea
% and obtain the estimated residuals $\hat{r}_i = r(\vy_i, \vxx^\star )$. This step can be solved with certifiable optimality using global solvers.

% \item {\bf Dual}: Update the weights with estimated residuals $\hat{r}_i$:
% \bea \label{eq:weightUpdate}
% \vw^\star = \argmin_{w_i \in [0,1]} \sumAllPointsi \hat{r}_i^2 w_i + \Phi(w_i)
% \eea
% This step can be solved in closed form given the form of $\Phi(w_i)$.
% \end{enumerate}
% \end{proposition}{}
% Decreasing $\mu$ increases the non-convexity (see Fig.~\ref{fig:GNCplots}(a)). The outlier process for $\rho_{\mu}(x)$ is:
% \bea \label{eq:outlierProcessGNCGM}
% \Phi(w_i) = \mu(\sqrt{w_i} - 1)^2
% \eea
% and the dual weight update rule is:
% \bea \label{eq:dualWeightUpdateGNCGM}
% w_i^\star = \left( \frac{\mu}{\hatr_i^2 + \mu} \right)^2
% \eea


% Increasing $\mu$ increases the non-convexity (see Fig.~\ref{fig:GNCplots}(b)). The outlier process for $\rho_{\mu}(x)$ is:
% \bea \label{eq:outlierProcessGNCTLS}
% \Phi(w_i) = \frac{\mu (1-w_i)}{ \mu + w_i} \barcsq
% \eea
% and the dual weight update rule is:
% \bea \label{eq:dualWeightUpdateGNCTLS}
% w_i^\star = \begin{cases}
% 0 & \hatr_i^2 \in \left[ \frac{\mu+1}{\mu}\barcsq, +\infty \right] \\
% \frac{\barc}{\hatr_i}\sqrt{\mu(\mu+1)} - \mu & \hatr_i^2 \in \left[ \frac{\mu}{\mu+1}\barcsq ,\frac{\mu+1}{\mu}\barcsq \right] \\
% 1 & \hatr_i^2 \in \left[ 0,\frac{\mu}{\mu+1}\barcsq \right]
% \end{cases}
% \eea